Yes. Let’s unpack **every single step** from a **camera-space 3D point** to a **2D screen point**, and then show **exactly where the Jacobian enters**.

I’ll do it in the cleanest order:

1. start with a 3D point in camera coordinates
2. send it through the graphics pipeline step by step
3. define the full map (h)
4. explain the local linearization
5. show how covariance uses that

---

# 0. What we are trying to compute

Suppose you already have a Gaussian center in **camera coordinates**:

[
x_c =
\begin{bmatrix}
X_c\
Y_c\
Z_c
\end{bmatrix}
]

This is a 3D point relative to the camera.

We want to know:

* where this point lands on the screen
* and how a **small cloud of nearby points** around it gets turned into a 2D ellipse

That second part is where covariance comes in.

---

# 1. Camera coordinates

At this stage, the point is in the camera’s own coordinate system.

So:

* (X_c): left-right location relative to camera
* (Y_c): up-down location relative to camera
* (Z_c): depth / distance in front of camera

This point is still 3D.

---

# 2. Convert to homogeneous coordinates

In graphics, points are often written in homogeneous form:

[
\tilde x_c =
\begin{bmatrix}
X_c\
Y_c\
Z_c\
1
\end{bmatrix}
]

Why?
Because matrix multiplication with (4\times 4) matrices becomes convenient.

---

# 3. Apply the perspective projection matrix

Now we multiply by the projection matrix (P):

[
\tilde x_{clip} = P \tilde x_c
]

This gives the **clip-space** point:

[
\tilde x_{clip} =
\begin{bmatrix}
x_{clip}\
y_{clip}\
z_{clip}\
w_{clip}
\end{bmatrix}
]

This is still **not** the final screen coordinate.

It is an intermediate representation.

---

## 3A. What does the projection matrix look like?

For a symmetric perspective camera, one common form is:

[
P =
\begin{bmatrix}
\frac{1}{\tan(\theta_x/2)} & 0 & 0 & 0\
0 & \frac{1}{\tan(\theta_y/2)} & 0 & 0\
0 & 0 & a & b\
0 & 0 & -1 & 0
\end{bmatrix}
]

or equivalently, in pixel-style intrinsics language, people often think in terms of:

* focal lengths
* principal point
* perspective divide

The exact matrix form depends on convention, but the important thing is:

**this step prepares the point for the perspective divide.**

---

# 4. Perspective divide

This is the crucial step.

After projection matrix multiplication, you divide by (w_{clip}):

[
x_{ndc} = \frac{x_{clip}}{w_{clip}},\qquad
y_{ndc} = \frac{y_{clip}}{w_{clip}},\qquad
z_{ndc} = \frac{z_{clip}}{w_{clip}}
]

These are the coordinates in **NDC** (normalized device coordinates).

This is where the mapping becomes nonlinear, because of the division.

That is the key step you were asking about earlier.

---

# 5. NDC coordinates

Now the point is inside a normalized cube:

* typically (x_{ndc}\in[-1,1])
* (y_{ndc}\in[-1,1])

This is often called the **canonical view volume** or **NDC space**.

So after the perspective divide, the point is no longer in real camera units.
It is in normalized coordinates.

---

# 6. Viewport transform: NDC (\to) screen pixels

Now convert NDC to actual screen pixels.

If the image width is (W) and height is (H), then one common mapping is:

[
u = \frac{W}{2}(x_{ndc}+1)
]

[
v = \frac{H}{2}(1-y_{ndc})
]

Now ((u,v)) is the actual screen pixel position.

So the pipeline so far is:

[
x_c \to \tilde x_c \to \tilde x_{clip} \to x_{ndc} \to s
]

where

[
s =
\begin{bmatrix}
u\
v
\end{bmatrix}
]

---

# 7. The full map (h)

Now let’s bundle all those steps together.

Define one function:

[
s = h(x_c)
]

That means:

[
h:\mathbb R^3 \to \mathbb R^2
]

and it takes a 3D camera-space point (x_c) and returns the 2D screen point (s).

So when I said earlier:

[
s = h(x_c)
]

I meant:

**all the little graphics steps together, treated as one big function.**

That function (h) includes:

1. homogeneous conversion
2. projection matrix
3. perspective divide
4. NDC to screen conversion

---

# 8. Why (h) is not linear

Because inside (h), there is a division by (w_{clip}).

That is the same reason perspective projection was nonlinear before.

So (h) is not a simple matrix multiplication from (\mathbb R^3\to\mathbb R^2).

That is why we cannot just write:

[
\Sigma_{screen} = A \Sigma_c A^T
]

with one fixed (A) for all points.

---

# 9. Small offsets around the Gaussian center

Now suppose the Gaussian center is at

[
\mu_c
]

A nearby point in the Gaussian is:

[
\mu_c + \delta x
]

where (\delta x) is a tiny 3D offset.

This means we are no longer looking at just one point.
We are looking at a tiny neighborhood around the center.

---

# 10. Exact screen mapping of that nearby point

The exact screen location of that nearby point is:

[
h(\mu_c+\delta x)
]

But this exact function is nonlinear and messy.

So we approximate it.

---

# 11. Local linear approximation

Near the center (\mu_c), the function (h) behaves approximately like a linear map.

So:

[
h(\mu_c+\delta x)\approx h(\mu_c) + J_h,\delta x
]

This is the important formula.

Let’s decode every part.

### (h(\mu_c))

This is the screen position of the Gaussian center.

### (\delta x)

Tiny 3D offset around the center.

### (J_h)

The Jacobian of the full camera-to-screen map (h), evaluated at (\mu_c).

It is a (2\times 3) matrix.

### (J_h \delta x)

Approximate 2D screen offset caused by the tiny 3D offset.

So:

[
\delta s \approx J_h \delta x
]

where

[
\delta s = h(\mu_c+\delta x)-h(\mu_c)
]

---

# 12. What is the Jacobian (J_h)?

It is the matrix of partial derivatives:

[
J_h =
\begin{bmatrix}
\frac{\partial u}{\partial X_c} & \frac{\partial u}{\partial Y_c} & \frac{\partial u}{\partial Z_c}\
\frac{\partial v}{\partial X_c} & \frac{\partial v}{\partial Y_c} & \frac{\partial v}{\partial Z_c}
\end{bmatrix}
]

This tells you:

* how screen (u) changes if (X_c) changes a tiny bit
* how screen (u) changes if (Y_c) changes a tiny bit
* how screen (u) changes if (Z_c) changes a tiny bit
* same for (v)

So (J_h) is a **tiny-motion converter**:

[
\text{tiny 3D camera offset} \to \text{tiny 2D screen offset}
]

---

# 13. Why covariance now uses the standard rule

Once we write

[
\delta s \approx J_h \delta x
]

the nearby offsets are now approximately related by a **linear map**.

That is the key.

For a linear map

[
y = A x
]

covariance transforms as

[
\Sigma_y = A \Sigma_x A^T
]

Here:

* (x = \delta x)
* (y = \delta s)
* (A = J_h)

So:

[
\Sigma_{screen} \approx J_h \Sigma_c J_h^T
]

That is the covariance formula.

---

# 14. So where exactly are we in the pipeline when we do this?

We are at this stage:

## First:

world Gaussian (\to) camera Gaussian

That gives (\Sigma_c).

## Then:

camera Gaussian (\to) screen ellipse

This second step is nonlinear because of perspective, so we use the Jacobian of the full camera-to-screen map.

So:

[
\Sigma_c \to \Sigma_{screen}
]

using

[
\Sigma_{screen} \approx J_h \Sigma_c J_h^T
]

---

# 15. Full chain from world covariance

If the world-to-camera linear part is (W), then:

[
\Sigma_c = W \Sigma_w W^T
]

and then

[
\Sigma_{screen} \approx J_h \Sigma_c J_h^T
]

So:

[
\Sigma_{screen} \approx J_h W \Sigma_w W^T J_h^T
]

That is the full chain.

---

# 16. Where did the skipped steps go?

They are all inside (h).

That is what I was compressing earlier.

Instead of spelling out:

* projection matrix
* clip space
* divide by (w)
* NDC
* viewport

I wrapped them into one function:

[
h(x_c)
]

Then (J_h) is the Jacobian of that full map.

So none of the steps disappeared.
They were just bundled together.

---

# 17. Two equivalent ways to think about it

## Way A: Direct computer vision style

Use:

[
u = f_x \frac{X_c}{Z_c}+c_x,\qquad
v = f_y \frac{Y_c}{Z_c}+c_y
]

Then compute Jacobian directly from that.

## Way B: Full graphics pipeline style

Use:

1. projection matrix
2. clip space
3. perspective divide
4. NDC
5. viewport transform

Then compute the Jacobian of the whole chain.

Both give the same idea.

---

# 18. The exact process, no skipped steps

Here is the full process as a chain.

## Point pipeline

[
x_c
\in \mathbb R^3
]

### Step 1:

homogeneous lift
[
x_c \to \tilde x_c =
\begin{bmatrix}
X_c\Y_c\Z_c\1
\end{bmatrix}
]

### Step 2:

projection matrix
[
\tilde x_{clip} = P \tilde x_c
]

### Step 3:

perspective divide
[
x_{ndc} =
\begin{bmatrix}
x_{clip}/w_{clip}\
y_{clip}/w_{clip}\
z_{clip}/w_{clip}
\end{bmatrix}
]

### Step 4:

viewport transform
[
s =
\begin{bmatrix}
u\v
\end{bmatrix}
]

So:

[
s = h(x_c)
]

---

## Covariance pipeline

Start with camera-space covariance:

[
\Sigma_c
]

Take small 3D offsets (\delta x) around the mean.

The exact screen offset is:

[
\delta s = h(\mu_c+\delta x)-h(\mu_c)
]

Approximate it linearly:

[
\delta s \approx J_h \delta x
]

Then covariance becomes:

[
\Sigma_{screen} \approx J_h \Sigma_c J_h^T
]

If you include world-to-camera:

[
\Sigma_c = W\Sigma_w W^T
]

then:

[
\Sigma_{screen} \approx J_h W\Sigma_w W^T J_h^T
]

---

# 19. What the code is conceptually doing

Yes, the code is conceptually doing all these steps.

But in implementation, it may:

* combine several steps
* use a direct closed-form Jacobian
* avoid building every intermediate quantity explicitly

So the math is the full story, even if the code is more compact.

---

# 20. The single sentence summary

Nothing was skipped mathematically: the steps were bundled into one full camera-to-screen function (h), and the Jacobian (J_h) is the derivative of that full chain, used to turn the 3D Gaussian covariance into a 2D screen-space covariance.

If you want, next I can do the **same full chain with actual numbers at every single stage**:
camera point (\to) clip (\to) NDC (\to) screen, and then compute (J_h).
